\documentclass[12pt,a4paper]{article}

\usepackage[utf8]{inputenc}
\usepackage[T1]{fontenc}
\usepackage[brazil]{babel}
\usepackage{amsmath}
\usepackage{booktabs}
\usepackage{geometry}
\usepackage{array}
\usepackage{float}
\usepackage{longtable}

\geometry{
    left=3cm,
    right=2cm,
    top=3cm,
    bottom=2cm
}

\renewcommand{\thesection}{\arabic{section}.}
\renewcommand{\thesubsection}{\arabic{section}.\arabic{subsection}.}

\title{\textbf{Questão 1 -- Construção e Verificação do Conjunto de Dados\\
Dados de Compartilhamento de Bicicletas}}
\author{}
\date{}

\begin{document}

\maketitle

\section{Introdução}

A Questão 1 tem como finalidade construir o conjunto de dados específico do grupo
a partir da base original \texttt{HW1\_bike\_sharing.csv} e verificar,
por meio de cálculos manuais e do ambiente R, os procedimentos estatísticos que
serão utilizados nas etapas posteriores do trabalho.

A base original contém 731 observações diárias referentes ao funcionamento de
um sistema de compartilhamento de bicicletas em uma cidade dos Estados Unidos.
Conforme definido no enunciado, cada grupo deve trabalhar com uma amostra
específica de 300 observações consecutivas, determinada a partir do maior número
de matrícula entre os integrantes.

Para o grupo, o valor obtido para o índice inicial foi $r=7$. Portanto, a
amostra utilizada é formada pelas observações de índice 7 até 306, totalizando
exatamente 300 registros consecutivos e preservando a ordem original dos dados.

A primeira observação da amostra corresponde ao dia 07 de janeiro de 2011,
enquanto a última corresponde ao dia 02 de novembro de 2011.

\section{Definição da amostra específica do grupo}

O índice inicial da amostra é definido no enunciado por

\begin{equation}
r = 1 + (M \bmod 100),
\end{equation}

em que $M$ representa o maior número de matrícula entre os integrantes do grupo.

A partir da definição da amostra utilizada no trabalho, obteve-se

\begin{equation}
r = 7.
\end{equation}

Assim, o resto da divisão do maior número de matrícula por 100 é igual a 6,
pois

\begin{equation}
7 = 1 + (M \bmod 100)
\end{equation}

e, consequentemente,

\begin{equation}
M \bmod 100 = 6.
\end{equation}

A seleção foi realizada mantendo as observações em sua ordem original, sem
exclusões ou reordenações. O intervalo utilizado pode ser resumido da seguinte
forma:

\begin{table}[H]
\centering
\caption{Identificação da amostra específica do grupo}
\begin{tabular}{lc}
\toprule
\textbf{Informação} & \textbf{Resultado} \\
\midrule
Índice inicial $r$ & 7 \\
Primeiro índice selecionado & 7 \\
Último índice selecionado & 306 \\
Número total de observações & 300 \\
Primeira data & 07/01/2011 \\
Última data & 02/11/2011 \\
\bottomrule
\end{tabular}
\end{table}

Dessa maneira, o conjunto selecionado atende à exigência de utilizar 300
observações consecutivas e constitui a base empregada nas análises subsequentes
do homework.

\section{Estruturação do conjunto de dados no R}

O conjunto selecionado contém as variáveis \texttt{instant}, \texttt{dteday},
\texttt{season}, \texttt{weathersit}, \texttt{temp}, \texttt{casual} e
\texttt{registered}.

No arquivo produzido pelo grupo, os dados são inicialmente organizados em uma
estrutura tabular no R. Em seguida, a coluna utilizada apenas como identificação
das linhas é transferida para os nomes das linhas e removida do corpo principal
da tabela. Também são definidos explicitamente os tipos de cada variável,
garantindo que datas, valores inteiros e valores numéricos sejam tratados de
forma adequada.

O código utilizado no arquivo do grupo é apresentado a seguir.

\begin{verbatim}
dados <- read.csv(
  text = texto_dados,
  stringsAsFactors = FALSE,
  check.names = FALSE
)

rownames(dados) <- dados$row
dados$row <- NULL

dados$dteday <- as.Date(dados$dteday)
dados$instant <- as.integer(dados$instant)
dados$season <- as.integer(dados$season)
dados$weathersit <- as.integer(dados$weathersit)
dados$temp <- as.numeric(dados$temp)
dados$casual <- as.integer(dados$casual)
dados$registered <- as.integer(dados$registered)

stopifnot(nrow(dados) == 300)
stopifnot(rownames(dados)[1] == "7")
stopifnot(rownames(dados)[300] == "306")

View(dados)

print(head(dados))
cat("\nTotal de observações:", nrow(dados), "\n")

library(openxlsx)

write.xlsx(
  dados,
  "data_group.xlsx",
  rowNames = TRUE,
  overwrite = TRUE
)
\end{verbatim}

As verificações realizadas com \texttt{stopifnot()} confirmam três condições
fundamentais: o conjunto possui exatamente 300 observações, a primeira linha
corresponde ao índice 7 e a última ao índice 306.

Ao final, o pacote \texttt{openxlsx} é utilizado para exportar a tabela para o
arquivo \texttt{data\_group.xlsx}, preservando os identificadores das linhas.

\section{Primeiras 10 observações da amostra}

Conforme solicitado no enunciado, os cálculos estatísticos foram realizados
inicialmente de forma manual utilizando somente as 10 primeiras observações do
conjunto selecionado.

Para essa etapa, foi considerada a variável \texttt{total\_user}, definida pela
soma do número de usuários casuais e registrados:

\begin{equation}
\text{total\_user}_i =
\text{casual}_i + \text{registered}_i.
\end{equation}

Os valores utilizados nos cálculos são apresentados na Tabela 2.

\begin{table}[H]
\centering
\caption{Primeiras 10 observações utilizadas nos cálculos manuais}
\begin{tabular}{cccccc}
\toprule
\textbf{Obs.} &
\textbf{instant} &
\textbf{Data} &
\textbf{casual} &
\textbf{registered} &
\textbf{total\_user} \\
\midrule
1  & 7  & 2011-01-07 & 148 & 1362 & 1510 \\
2  & 8  & 2011-01-08 & 68  & 891  & 959  \\
3  & 9  & 2011-01-09 & 54  & 768  & 822  \\
4  & 10 & 2011-01-10 & 41  & 1280 & 1321 \\
5  & 11 & 2011-01-11 & 43  & 1220 & 1263 \\
6  & 12 & 2011-01-12 & 25  & 1137 & 1162 \\
7  & 13 & 2011-01-13 & 38  & 1368 & 1406 \\
8  & 14 & 2011-01-14 & 54  & 1367 & 1421 \\
9  & 15 & 2011-01-15 & 222 & 1026 & 1248 \\
10 & 16 & 2011-01-16 & 251 & 953  & 1204 \\
\bottomrule
\end{tabular}
\end{table}

Esses 10 valores de \texttt{total\_user} foram utilizados para calcular média,
mediana, moda, variância amostral, desvio padrão amostral, quartis e intervalo
interquartil.

\section{Cálculos manuais}

\subsection{Média aritmética}

A média aritmética é calculada pela soma de todos os valores observados dividida
pelo número total de observações:

\begin{equation}
\bar{x} =
\frac{\sum_{i=1}^{n} x_i}{n}.
\end{equation}

Para as 10 primeiras observações,

\begin{equation}
\bar{x} =
\frac{
1510 + 959 + 822 + 1321 + 1263 +
1162 + 1406 + 1421 + 1248 + 1204
}{10}.
\end{equation}

A soma dos valores é

\begin{equation}
\sum_{i=1}^{10}x_i = 12316.
\end{equation}

Portanto,

\begin{equation}
\bar{x} =
\frac{12316}{10}
=
1231{,}6.
\end{equation}

Assim, a média do número total de usuários nas 10 primeiras observações é
1231,6 usuários.

\subsection{Mediana}

Para determinar a mediana, os valores de \texttt{total\_user} são colocados em
ordem crescente:

\begin{center}
822,\ 959,\ 1162,\ 1204,\ 1248,\ 1263,\ 1321,\ 1406,\ 1421,\ 1510.
\end{center}

Como existem 10 observações, o número de valores é par. Nesse caso, a mediana é
obtida pela média aritmética dos dois valores centrais, localizados nas posições
5 e 6:

\begin{equation}
Md =
\frac{x_{(5)} + x_{(6)}}{2}.
\end{equation}

Substituindo os valores,

\begin{equation}
Md =
\frac{1248 + 1263}{2}
=
\frac{2511}{2}
=
1255{,}5.
\end{equation}

Portanto, a mediana é igual a 1255,5 usuários.

\subsection{Moda}

A moda corresponde ao valor que apresenta a maior frequência de ocorrência.

Nas 10 observações analisadas, todos os valores de \texttt{total\_user}
aparecem apenas uma vez. Como não existe um valor com frequência superior aos
demais, o conjunto é classificado como amodal.

\subsection{Variância amostral}

A variância amostral é calculada pela soma dos quadrados das diferenças entre
cada observação e a média, dividida por $n-1$:

\begin{equation}
s^2 =
\frac{
\sum_{i=1}^{n}(x_i-\bar{x})^2
}{n-1}.
\end{equation}

Utilizando $\bar{x}=1231{,}6$, obtêm-se os valores apresentados na Tabela 3.

\begin{table}[H]
\centering
\caption{Cálculo das diferenças ao quadrado para a variância amostral}
\begin{tabular}{rrrr}
\toprule
\textbf{Obs.} &
\textbf{$x_i$} &
\textbf{$x_i-\bar{x}$} &
\textbf{$(x_i-\bar{x})^2$} \\
\midrule
1  & 1510 & 278,4  & 77506,56 \\
2  & 959  & -272,6 & 74310,76 \\
3  & 822  & -409,6 & 167772,16 \\
4  & 1321 & 89,4   & 7992,36 \\
5  & 1263 & 31,4   & 985,96 \\
6  & 1162 & -69,6  & 4844,16 \\
7  & 1406 & 174,4  & 30415,36 \\
8  & 1421 & 189,4  & 35872,36 \\
9  & 1248 & 16,4   & 268,96 \\
10 & 1204 & -27,6  & 761,76 \\
\midrule
\textbf{Soma} & \textbf{12316} & \textbf{0} & \textbf{400730,40} \\
\bottomrule
\end{tabular}
\end{table}

Logo,

\begin{equation}
s^2 =
\frac{400730{,}40}{10-1}
=
\frac{400730{,}40}{9}
=
44525{,}60.
\end{equation}

A variância amostral encontrada é, portanto, 44525,60.

\subsection{Desvio padrão amostral}

O desvio padrão amostral é obtido pela raiz quadrada da variância:

\begin{equation}
s = \sqrt{s^2}.
\end{equation}

Substituindo o valor calculado,

\begin{equation}
s =
\sqrt{44525{,}60}
\approx
211{,}01.
\end{equation}

Assim, o desvio padrão amostral das 10 primeiras observações é aproximadamente
211,01 usuários.

\subsection{Quartis e intervalo interquartil}

No cálculo manual realizado pelo grupo, os quartis foram determinados por meio
de interpolação utilizando posições baseadas em $(n+1)$.

Para o primeiro quartil,

\begin{equation}
\text{posição de }Q_1 =
\frac{n+1}{4}
=
\frac{11}{4}
=
2{,}75.
\end{equation}

A posição 2,75 encontra-se entre o segundo valor, 959, e o terceiro valor, 1162.
Assim,

\begin{equation}
Q_1 =
959 + 0{,}75(1162-959)
=
1111{,}25.
\end{equation}

O segundo quartil corresponde à mediana:

\begin{equation}
Q_2 = 1255{,}5.
\end{equation}

Para o terceiro quartil,

\begin{equation}
\text{posição de }Q_3 =
\frac{3(n+1)}{4}
=
\frac{33}{4}
=
8{,}25.
\end{equation}

A posição 8,25 encontra-se entre o oitavo valor, 1406, e o nono valor, 1421:

\begin{equation}
Q_3 =
1406 + 0{,}25(1421-1406)
=
1409{,}75.
\end{equation}

O intervalo interquartil é definido por

\begin{equation}
IQR = Q_3-Q_1.
\end{equation}

Portanto,

\begin{equation}
IQR =
1409{,}75-1111{,}25
=
298{,}50.
\end{equation}

\section{Verificação dos cálculos no ambiente R}

Após a realização manual, os cálculos foram verificados no ambiente R utilizando
as mesmas 10 observações.

O procedimento utilizado pelo grupo considera as funções \texttt{mean()},
\texttt{median()}, \texttt{var()}, \texttt{sd()}, \texttt{quantile()} e
\texttt{IQR()} para reproduzir computacionalmente as medidas calculadas.

Os resultados obtidos no R foram:

\begin{table}[H]
\centering
\caption{Resultados obtidos no R para as 10 primeiras observações}
\begin{tabular}{lc}
\toprule
\textbf{Medida} & \textbf{Resultado} \\
\midrule
Média & 1231,60 \\
Mediana & 1255,50 \\
Variância amostral & 44525,60 \\
Desvio padrão amostral & 211,0109 \\
$Q_1$ & 1172,50 \\
$Q_2$ & 1255,50 \\
$Q_3$ & 1384,75 \\
IQR & 212,25 \\
\bottomrule
\end{tabular}
\end{table}

\section{Comparação entre os cálculos manuais e os resultados do R}

A comparação entre os resultados permite verificar quais medidas coincidiram
entre o procedimento manual e o procedimento computacional.

\begin{table}[H]
\centering
\caption{Comparação entre os cálculos manuais e os resultados do R}
\begin{tabular}{lccc}
\toprule
\textbf{Estatística} &
\textbf{Cálculo manual} &
\textbf{R} &
\textbf{Resultado} \\
\midrule
Média & 1231,60 & 1231,60 & Idêntico \\
Mediana & 1255,50 & 1255,50 & Idêntico \\
Variância amostral & 44525,60 & 44525,60 & Idêntico \\
Desvio padrão amostral & 211,01 & 211,01 & Idêntico \\
$Q_1$ & 1111,25 & 1172,50 & Diferente \\
$Q_3$ & 1409,75 & 1384,75 & Diferente \\
IQR & 298,50 & 212,25 & Diferente \\
\bottomrule
\end{tabular}
\end{table}

A média, a mediana, a variância amostral e o desvio padrão amostral apresentaram
os mesmos resultados nos dois procedimentos. Dessa forma, os cálculos manuais
dessas medidas foram confirmados pelas funções correspondentes do R.

As diferenças foram observadas em $Q_1$, $Q_3$ e, consequentemente, no intervalo
interquartil. Conforme registrado nos cálculos do grupo, essa divergência ocorre
porque o procedimento manual e a função \texttt{quantile()} do R utilizam
convenções diferentes para determinar percentis.

No cálculo manual foi empregada uma posição baseada em $(n+1)$ e interpolação
linear entre os valores adjacentes. Já a função \texttt{quantile()} do R utiliza,
por padrão, o método denominado \texttt{type = 7}, que adota outra regra de
interpolação.

Por esse motivo, a diferença entre os quartis não representa um erro aritmético,
mas uma consequência direta do método utilizado. Para as etapas seguintes do
trabalho, foram adotados os valores calculados pelo R, mantendo um único padrão
computacional para a análise das 300 observações.

\section{Aplicação ao conjunto completo}

Depois de verificar os procedimentos com as 10 primeiras observações, os mesmos
métodos estatísticos foram estendidos ao conjunto completo de 300 registros.

Essa etapa permite utilizar uma base previamente conferida nos itens seguintes
do homework. O conjunto específico do grupo permanece formado pelas observações
7 a 306, correspondentes ao período de 07/01/2011 a 02/11/2011.

A verificação manual cumpre, portanto, a função de conferir os procedimentos
estatísticos em uma amostra pequena e controlável antes da aplicação dos mesmos
comandos ao conjunto completo.

\section{Conclusão}

A Questão 1 estabeleceu a base de dados específica que será utilizada nas
demais etapas do trabalho. A partir do critério definido no enunciado, foi
obtido $r=7$, resultando na seleção de 300 observações consecutivas, dos índices
7 a 306, correspondentes ao período entre 07 de janeiro de 2011 e 02 de novembro
de 2011.

Nas 10 primeiras observações, a variável \texttt{total\_user} apresentou média
igual a 1231,6 e mediana igual a 1255,5, enquanto a variância amostral foi
44525,60 e o desvio padrão amostral aproximadamente 211,01. Esses resultados
coincidiram com os valores calculados no R.

Para os quartis, foram obtidos valores diferentes entre o cálculo manual e o R.
O cálculo manual resultou em $Q_1=1111{,}25$, $Q_3=1409{,}75$ e
$IQR=298{,}50$, enquanto o R retornou $Q_1=1172{,}50$,
$Q_3=1384{,}75$ e $IQR=212{,}25$. A diferença decorre das distintas convenções
de interpolação utilizadas para o cálculo dos percentis.

Com a conferência dos procedimentos, o conjunto completo de 300 observações
fica preparado para servir de base às análises estatísticas desenvolvidas nas
questões seguintes.

\end{document}
